\section{Completed primary and random-feature experiments}
\label{sec:primary-results}
All 204 registered runs completed through 60200 updates, and the independent raw-record audit passed for all runs. This section reports the 60 main and 24 random-feature runs; Section~\ref{sec:scope} reports all 120 recipe/split/operation controls. Figure~\ref{fig:width} and Table~\ref{tab:main} retain every main width and denominator; supporting trajectories and every seed clock appear in Appendix~\ref{app:primary-displays}.

\begin{figure}[t]
\centering
\includegraphics[width=\linewidth]{width_summary.pdf}
\caption{Main-recipe width comparison, all twelve blocks. (a) Seed held-out onsets, restricted train/held-out means and 95\% seed-bootstrap intervals. Open arrows mark non-events restricted at $H=60000$, rather than observed crossings. (b) Generalization and delayed-event fractions with Wilson 95\% intervals. (c) All paired width-512-minus-32 restricted generalization-time differences and their mean/95\% paired interval. Collapsed intervals under common censoring do not establish eventual equality.}
\label{fig:width}
\end{figure}
\begin{table}[t]
\caption{All main widths: instantiated total/trainable parameters, training-fit, generalization and delayed-event counts, and mean restricted generalization time $R_g$ in \emph{thousand updates} with 95\% seed-bootstrap intervals. All twelve blocks remain in every denominator; horizon values are restrictions, not observed onsets.}
\label{tab:main}
\centering
\input{../tables/main_outcomes}
\end{table}

\subsection{Width changes held-out learning within the window}
All twelve main runs at each width reach training fit. Confirmed generalization counts at widths $32,64,128,256,512$ are respectively $0,0,8,12,12$ out of twelve; delayed-event counts are the same. Thus fitting is insufficient for generalization by the common horizon, and the observed occurrence of delayed generalization increases across this width grid. Width~128 retains four noncrossings rather than an imputed eventual onset. Onset and terminal sustained-generalization counts agree in these completed cells, without implying continuous success between evaluations or beyond execution.

For P1, the paired width~512-minus-32 restricted generalization-time difference is $-43800$ updates (95\% paired-bootstrap interval $[-44150,-43433]$, twelve blocks). This supports its prespecified negative direction. Width~512's mean observed onset is 16200 updates ($[15850,16567]$); width~32 contributes restriction at 60000 in every block. This is a finite-window contrast; width~32's degenerate restricted interval does not locate its eventual onset.

Training fit is also earlier at larger width: the endpoint paired restricted-fit difference is $-1742$ updates ($[-1817,-1675]$). The paired restricted waiting-time difference remains $-42058$ updates ($[-42408,-41717]$), so the observed generalization-clock contrast cannot be described solely as moving the fit clock earlier. This does not establish that eventual post-fit delay is shorter for all narrow models, whose generalization times remain censored.

\subsection{Substantial delay persists at width~512}
P2 is supported: all twelve width~512 blocks meet the registered lag and ratio criteria, giving delayed-event prevalence $12/12$ with Wilson 95\% interval $[0.758,1]$. The lower bound exceeds the declared majority boundary. Every run fits at update 200 on the sampled grid, while the mean post-fit wait is 16000 updates ($[15650,16367]$). These data show earlier generalization alongside a substantial remaining delay. Training-fit accuracy does not identify an internal memorization circuit.

\subsection{The frozen control does not resolve the explanation}
Figure~\ref{fig:frozen} retains the matched intervention display in the appendix. In the first eight matched blocks, the full-minus-frozen endpoint interaction is $-43975$ updates ($[-44325,-43688]$). Its statistical direction agrees with P3, but its prospective fit gate fails. The frozen width~32 models fit in $0/8$ runs and frozen width~512 in $8/8$, compared with $8/8$ at each full-model endpoint; frozen width~128 also fits in $0/8$. No frozen width has a confirmed generalization event in any of its eight runs.

Thus the narrow frozen model cannot supply the planned fit-comparable discriminator. At width~512 the frozen model can fit yet does not generalize by the horizon, providing a within-width observation consistent with the importance of hidden trainability under this recipe. It does not rescue the failed across-width gate or distinguish adaptation from hidden decay, optimization geometry and trainable count. The complete candidate explanation is therefore \emph{not established}: P1 and P2 are supported, while the representation-versus-fitting interpretation of P3 remains unresolved.

\subsection{Definition sensitivity and the compute estimand}
All thirteen prescribed sensitivity variants preserve the negative P1 interval, P2's majority decision and P3's failed fit gate on these completed families. The held-out threshold changes the measured clock contrast; it does not select a new primary definition. The 200-update thinned view only removes existing evaluations and has reduced persistence-confirmation coverage near the endpoint. Scope-family sensitivities are reported below and in Appendix~\ref{app:scope-displays}.

The update-based ranking differs from a fixed measured CPU budget. At the last measured evaluation at or before 100 process-CPU seconds, width~128's mean held-out accuracy is 92.26\% ($[91.18,93.26]$) versus 61.17\% ($[58.25,64.02]$) at width~512, with twelve available blocks in each cell. The actual sampled steps span 27700--31200 and 9700--10100, respectively; no interpolation to an exact 100-second value is performed. These are descriptive all-block summaries on this implementation and host, rather than a paired test of an optimal width or a hardware-independent efficiency law.
